% !PRINTED 517
\weberpage{517}
%<<CONT
are. If, then, we denote by $p, p', p'', \ldots$ the prime factors of $m$ of the form $4n+1$, and by $q, q', q'', \ldots$ the prime factors of $m$ of the form $4n+3$, then the following square roots occur in the partial equations:
\[
\sqrt{p}, \sqrt{p'}, \sqrt{p''}, \ldots, \quad \sqrt{qq'}, \sqrt{qq''}, \ldots;
\]
so that if only one $q$ is contained in $m$, then $\sqrt{q}$ does not occur in the partial equations. Since in this case at least one $q$ must divide $m$, the number of square roots required for the decomposition is equal to the number of primes dividing $m$, diminished by $1$, and hence equal to the number of independent characters corresponding to this case.

2) If
\[
m\equiv 1\ (\mathrm{mod}\ 4), \quad m\equiv 6, 2\ (\mathrm{mod}\ 8), \quad m\equiv 4\ (\mathrm{mod}\ 16),
\]
then, by § 105, the characters
\[
\left(\frac{-1}{A}\right), \ \left(\frac{2}{A}\right), \ \left(\frac{-2}{A}\right), \ \left(\frac{-1}{A}\right)
\]
can respectively be expressed by the characters of the form
\[
\left(\frac{A}{m''}\right)
\]
and so in these cases too formula (8) suffices to separate all the genera from one another.

Formulae (9), (12) teach that in these cases the following adjunctions are necessary for the complete separation of the genera:
\weq{17}{
\begin{aligned}
m&\equiv 1\ (\mathrm{mod}\ 4), & \sqrt{p}, \sqrt{p'}, \sqrt{p''}, \ldots, \ & \sqrt{q}, \sqrt{q'}, \sqrt{q''}, \ldots \\
m&\equiv 6\ (\mathrm{mod}\ 8), & \sqrt{2}, \sqrt{p}, \sqrt{p'}, \sqrt{p''}, \ldots, \ & \sqrt{qq'}, \sqrt{qq''}, \ldots \\
m&\equiv 2\ (\mathrm{mod}\ 8), & \sqrt{p}, \sqrt{p'}, \sqrt{p''}, \ldots, \ & \sqrt{2q}, \sqrt{2q'}, \sqrt{2q''}, \ldots \\
m&\equiv 4\ (\mathrm{mod}\ 16), & \sqrt{p}, \sqrt{p'}, \sqrt{p''}, \ldots, \ & \sqrt{q}, \sqrt{q'}, \sqrt{q''}, \ldots
\end{aligned}
}

For in the case $m\equiv 1 \ (\mathrm{mod}\ 4)$ the number of the $q$ is in any event even. If, then, we put $m = q r^2 m''$, $m' = q r^2$ and understand by $r^2$ the greatest square number dividing $r^2 m''$, then $m''\equiv 3 \ (\mathrm{mod}\ 4)$, and formula (12) gives the adjunction of $\sqrt{q}$; thus in this case, for the perfect separation of the genera,
\[
\sqrt{p}, \sqrt{p'}, \sqrt{p''}, \ldots, \quad \sqrt{q}, \sqrt{q'}, \sqrt{q''}, \ldots
\]
are to be adjoined, and the number of square roots is equal to the number of primes dividing $m$, again in agree-
%CONT>>
